\documentclass[11pt]{article}
\usepackage[margin=1in]{geometry}
\usepackage{amsmath,amssymb,amsfonts,amsthm,mathtools,bm}
\usepackage[T1]{fontenc}
\usepackage{microtype}
\newtheorem{theorem}{Theorem}
\newtheorem{lemma}{Lemma}
\newtheorem{cor}{Corollary}
\newtheorem{defn}{Definition}
\newtheorem{principle}{Principle}
\newtheorem{statement}{Statement}
\theoremstyle{remark}
\newtheorem{statementr}{Statement}

\title{Volume Vectors, Determinants, Trace, and Characteristic Equations}
\author{James Taylor}
\date{Working draft derived from the 2003 thesis}
\begin{document}
\maketitle

\section{Supplement to Chapter \ref{c:algo}} \label{a:app7}


We start with some useful notations. 
We denote the set of the first $n$ numbers of the
natural numbers, i.e.\ $\{1, \ldots n\}$, by \ensuremath{\mathfrak{n}}.
We use \ensuremath{\ensuremath{\mathfrak{n}}^{k}_{\neq}}\ for the set of
$k$-tuples, each entry being distinct, with the entries being
elements of \ensuremath{\mathfrak{n}} while  \ensuremath{{}^{k}\!\ensuremath{\mathfrak{n}}} is the set of $k$-element subsets
of \ensuremath{\mathfrak{n}}. We shall use the notation $(I)$, where $I$ is a subset of
a canonically ordered set (such as the natural numbers), to indicate
we wish to form an ordered set using the canonical ordering,
e.g. $(\{1, 5,  3\}) \ensuremath{:=} (1,3,5)$. Conversely, if we have an ordered
set $I$, then we shall use $\{ I \}$ to indicate the set itself. Thus,
$\{ \}$ is the inverse of $( )$. In this section, they will be maps
between \ensuremath{\ensuremath{\mathfrak{n}}^{k}_{\neq}}\ and \ensuremath{{}^{k}\!\ensuremath{\mathfrak{n}}}. All tensor and wedge products
are utilizing the idea of the set-based tensor product of appendix
\ref{a:tensor}. In particular, the ordering in which things are
written down is not relevant; only the labelling is relevant. 


We also define an anti-symmetric symbol which is crucial for component
formulas. The symbol $\ensuremath{\epsilon^{n}_{I}}$ is defined to be the sign of the
permutation putting the numbers in the correct order, i.e.\ taking 
$((\ensuremath{\mathfrak{n}}\setminus \{I\}), I) \to (\ensuremath{\mathfrak{n}})$ or 
$((\ensuremath{\mathfrak{n}}\setminus I), (I)) \to (\ensuremath{\mathfrak{n}})$ depending on whether $I$
is an ordered set or unordered. Just to emphasize, we will be using
this symbol both for ordered and unordered sets.


Let $A:V \to W$ be a map between two vector spaces of dimension
$n$. Each vector space is endowed with a
volume vector, say  $\nu$ for $V$ and $\omega$ for $W$. By a volume
vector, say for $V$, we mean a non-zero  element of $\bigwedge^n V$, which is
the 1-dimensional vector space formed from the totally
antisymmetric subspace of the $n$-fold self-tensor product of $V$.
By $A^{\otimes n}$, we mean the linear operator acting on $V^{\otimes
n}$ 
taking
values in $W\ensuremath{^{\otimes n} }$ defined on product elements by
\[
A\ensuremath{^{\otimes n} } (\mathop{\otimes}_{i \in \ensuremath{\mathfrak{n}}} e_i ) = \mathop{\otimes}_{i \in \ensuremath{\mathfrak{n}}} A(e_i).
\] 
A key property of $A\ensuremath{^{\otimes n} }$ is that it preserves spaces invariant
under permutations. That is to say, if $P_{\ensuremath{\sigma}}$ represents the
action of a 
permutation \ensuremath{\sigma}\  on
an $n$-fold self-tensor product induced by the permutation acting on
the indices, then $A \ensuremath{^{\otimes n} } P_{\ensuremath{\sigma}} = P_{\ensuremath{\sigma}} A \ensuremath{^{\otimes n} }$. 
Thus, $A \ensuremath{^{\otimes n} } (\nu) \in
\bigwedge^n V$ since  $P_{\ensuremath{\sigma}} A \ensuremath{^{\otimes n} }
(\nu)
= -1^{\ensuremath{\sigma}} A \ensuremath{^{\otimes n} } (\nu)$ for every permutation $\ensuremath{\sigma}$. This
characterizes a volume vector. 

We define the determinant of $A$ as the unique number satisfying
\begin{equation} \label{7e:detdef}
A\ensuremath{^{\otimes n} } (\nu) \ensuremath{=:} (\det A) \omega.
\end{equation} 
If we choose 
bases $\{v_i\}$ and $\{w_i\}$ such that the volume vectors are the wedge
product of the bases in the given order, i.e.\ $\nu = \bigwedge_{i \in
\ensuremath{\mathfrak{n}}} v_i$, then the
 determinant is the number satisfying 
\begin{equation} \label{ae:det}
\bigwedge_{i\in \ensuremath{\mathfrak{n}}} A( v_i) \ensuremath{=:} (\det A)  \bigwedge_{i
\in\ensuremath{\mathfrak{n}}} 
w_i.
\end{equation}
In its turn, we find the usual determinant formula 
\[
\det A = \sum_{\ensuremath{\sigma}: \ensuremath{\mathfrak{n}} \leftrightarrow \ensuremath{\mathfrak{n}}} -1^{\ensuremath{\sigma}}
\prod_{i \in n} A_i^{\ensuremath{\sigma}(i)}.  
\]
We also note that if $V=W$, then we can choose $\nu = \omega$. This
means that the determinant in that case is defined independently of the
choice of $\nu$ since different choices correspond to different
constants and these constants can be cancelled. 

The trace of an operator can also be defined using volume vectors and
tensors. But its definition requires that $A$ is map from $V$ to
itself. Then $A^{(i)}$ is defined to be 
the operator acting on $V\ensuremath{^{\otimes n} }$ defined
on product elements via
\[
A^{(i)}  (\mathop{\otimes}_{j \in \ensuremath{\mathfrak{n}}} e_j) = \{\mathop{\otimes}_{j \in \ensuremath{\mathfrak{n}}
\setminus {\{i\}} } e_j \} \otimes A(e_i).
\]
That is to say, it is the identity on each factor except the \ensuremath{{i}^{th}}\
factor for which $A$ acts on that factor.
The operator $\sum_{i \in \ensuremath{\mathfrak{n}}} A^{(i)}$  clearly commutes with
permutations. Thus, one can define the trace of the operator $A$ as
the unique number $\ensuremath{\textrm{Tr}} A$ satisfying 
\[
\sum_{i \in \ensuremath{\mathfrak{n}}} A^{(i)} (\nu) \ensuremath{=:} (\ensuremath{\textrm{Tr}} A) \nu.
\] 
In coordinates, we see that 
\[
\sum_{i \in \ensuremath{\mathfrak{n}}} A^{(i)} (\nu) = \sum_{i \in \ensuremath{\mathfrak{n}}} \{\bigwedge_{j
\in j \in \ensuremath{\mathfrak{n}}
\setminus {\{i\}}} v_j \} \wedge A(v_i).
\]
Since $A(v_i) = \sum_k A_i^k v_k$, we notice that in the \ensuremath{{i}^{th}} term of
the sum, only $A_i^i v_i$ gives a non-zero contribution to the wedge
product. Thus, we have that the trace is given, in components, by 
\[
\ensuremath{\textrm{Tr}} A = \sum_i A_i^i
\]
which is just the usual coordinate formula. Also notice that this
definition is independent of the volume form chosen and hence any
wedge product of a basis will yield the same result. In other words,
the trace is canonically defined.\footnote{The reason that we need $A$ to map
$V$ to itself, rather than some other vector space, is that the
identity is crucial in this argument. On the other hand, if $A: V \to
W$, then if we choose the volume vectors $\nu$ and $\omega$ as well as
another map $B: V \to W$, then we can define 
\[
\ensuremath{\textrm{Tr}}_B A (\omega) = (\sum_{i \in \ensuremath{\mathfrak{n}}} \{\mathop{\otimes}_{j \in \{\ensuremath{\mathfrak{n}}
\setminus i\} }B \} \otimes (A)_i) \nu
\]
Taking $A =B$, we find that the above formula gives us $n$ times the
determinant of $A$.  One can also envision a further generalization
involving multiple operators appearing a specified number of
times. This begins to resemble something remarkably similar to minors.
}

We needed to choose a volume form on each space in order 
to define the determinant. One way to do this is to specify a
metric. Then, up to a sign, we have a canonical volume vector. This is
formed by taking an orthonormal basis and wedging all of the elements
in the basis together. The sign ambiguity comes in by not knowing how
to choose which order to wedge product the vectors together. But, if
we want a positive number, then there is no ambiguity. In fact, given
$A$, we can redefine one of the volume vectors so that the determinant
is positive. Thus, $\ensuremath{\lvert {\det A} \rvert}$ or just $\ensuremath{\lvert {A} \rvert}$ for short, is
well-defined. 
Note that if we had a map between $V$ and itself, then given
any volume vector on $V$, we would use the same volume vector on both
sides. This gives an unambiguous determinant which may be negative.  

Extending this to manifolds should now be clear. 
Our setup is that we have a mapping $f: M^k \to N^k$ between two
Riemannian manifolds of the same dimension. 
Since $D_pf : T_pM \to T_{f(p)}N$, we can ask what is the
determinant of $Df$ as a mapping of the tangent spaces. The volume
vectors that we choose are the Riemannian volume vectors. For $M$, let
$\nu$ be the volume form and $\omega$ be the volume form for $N$. Then
$\ensuremath{\lvert {\det(D_pf)} \rvert} \ensuremath{=:} \ensuremath{\lvert {D_pf} \rvert} $   satisfies
\[
\ensuremath{\lvert {D_pf} \rvert} \omega_{f(p)} =( D_pf )\ensuremath{^{\otimes k} } (\nu_p)
\]
up to a sign.
That
completely determines the information. In any coordinate system, the
volume form is $\sqrt{\det G^{ij}} \bigwedge_{i \in \ensuremath{\mathfrak{n}}} v_i$.  Thus, the
determinant in coordinates is
\[
\ensuremath{\lvert {D_pf} \rvert} = \ensuremath{\lvert {\frac{\sqrt{\det{G_M^{ij}}}}{\sqrt{\det{G_N^{ij}}}} 
\sum_{\sigma : \ensuremath{\mathfrak{n}} \leftrightarrow \ensuremath{\mathfrak{n}}} -1^\sigma
\prod_{j \in \ensuremath{\mathfrak{n}}} \ensuremath{\frac{\partial {f^{\sigma j} }}{\partial {x^j}}} } \rvert}
\]
We do not have to worry about orientability as we are concerned only
with local properties of objects which should be non-zero.\footnote{If 
we were concerned about the Maslov index and the
determinant vanishing, we would choose an initial volume vector in a
neighborhood about the initial point and
parallel transport the volume vector along the curve. As we pass through
a zero of the determinant, we would pickup a phase and appropriately
order the volume vector to have a positive determinant.}


\section{Derivatives of the determinant}

We need to compute derivatives of the volume vector. Our setup is that
we have a map $Q$ with differential $DQ$ acting on the tangent
space; we shall suppress $t$ in the following discussion. 
The only important fact is that the Riemannian 
volume vector is parallel. The reason is that we can choose an
orthonormal basis at a point $q \in \ensuremath{\mathcal{Q}}$ and then parallel
transport the basis along geodesics. This gives us a basis of
orthonormal vector fields at every point as the inner product is
preserved by parallel transport. 


We have that  the first covariant derivative of these
vector fields is zero at $q$. Thus, at every point in \ensuremath{\mathcal{Q}}, the Riemannian
volume form's first covariant derivative vanishes. This means that it
is parallel and that the higher derivatives also vanish since these
will be derivatives of a zero tensor.


With that trivial fact established, we can now differentiate the
determinant. Let $A \ensuremath{:=} DQ$. We start with 
\[
\ensuremath{\lvert {A} \rvert}  \omega  = A\ensuremath{^{\otimes n} } \circ \nu.
\]
Recall that composition essentially means that we tensor these objects
together and then contract the indices. Since covariant derivatives
commute  with contractions, we have that
\[
D[\ensuremath{\lvert {A} \rvert}]  \omega + \ensuremath{\lvert {A} \rvert}  D[\omega]
=
D[A\ensuremath{^{\otimes n} }] \circ \nu +  A\ensuremath{^{\otimes n} } \circ D[\nu].
\]
Since the volume forms are parallel, this reduces to 
\[
D[\ensuremath{\lvert {A} \rvert}] \omega = D[A\ensuremath{^{\otimes n} }] \circ \nu.
\]
Continuing in this fashion, we have
\[
D^m [\ensuremath{\lvert {A} \rvert}]  \omega = D^m[A\ensuremath{^{\otimes n} }] \circ \nu.
\]
One can then use the product rule to express $D^m[A\ensuremath{^{\otimes n} }]$ in terms
of the factors. If one does this and then expresses it in coordinates,
we find
\[
(D^{\ensuremath{\lvert {i} \rvert}} [\det(A)])_{k_1\cdots k_i} = 
\sum_{\ensuremath{PR} \in \ensuremath{\mathcal{\ensuremath{PR}}} (i, n)} 
 \sum_{\sigma : \ensuremath{\mathfrak{n}} \leftrightarrow \ensuremath{\mathfrak{n}}}
-1^{\sigma}
\prod_{j \in \ensuremath{\mathfrak{n}}} A^{\sigma_j}_{j k_{r_j(1)} \cdots k_{r_j(\ensuremath{\lvert {r_j} \rvert})}}
\]
where $\ensuremath{\mathcal{\ensuremath{PR}}}(i, n)$ is the set of 
 product rule diagrams of length $i$ with height $n$, $R_{\ensuremath{PR}}$ is the set of
nonempty rows, and   $r_j$ is the \ensuremath{{j}^{th}}\ row. The extra subscripts
 indicate covariant differentiation slots. This is the same
 formula as in the Euclidean case; it can also be written in terms of
 minors, although we shall not do that here.  


\section{The continuity equation}


The continuity equation is 
\[
\ensuremath{\frac{\partial {\ensuremath{\rho}}}{\partial {t}}} + \nabla \cdot (\ensuremath{\rho}  \ensuremath{\boldsymbol{v}}) =0
\]
where $\ensuremath{\rho} : \ensuremath{\mathcal{Q}} \to \ensuremath{\mathbb{R}^{}}$ and \ensuremath{\boldsymbol{v}}\ is a time-dependent
velocity field on \ensuremath{\mathcal{Q}}. We assume that $Q(x_0,t)$ describes the
family of integral curves of the velocity field. 
Let $\ensuremath{\rho}_0$ be given. Then we claim that the solution to the
continuity equation, as said in Chapter \ref{c:algo}, 
is
\[
\ensuremath{\rho}_t(x) = \ensuremath{\rho}_0(Q^{-1}(x, t))\ensuremath{\lvert {DQ^{-1} (x,t)} \rvert}.
\]
To understand this, we note that under a mapping $f$, a density $\mu$ on the
image is mapped to a density on the domain via
$ \tilde{\mu}(x) \ensuremath{:=} \mu(f(x))\ensuremath{\lvert {Df} \rvert}(x)$. This is the definition
which gives $\int_A \tilde{\mu} = \int_{f(A)} \mu$. The Jacobian
factor is what is needed for the change of variables. This is the
picture of why it is true.

Direct verification is as follows. 
Define $h(x,t) \ensuremath{:=}
Q^{-1}(x,t)$. Then $\ensuremath{\rho}(x,t) \ensuremath{:=} \ensuremath{\rho}_0(h(x,t)) \ensuremath{\lvert {Dh} \rvert} (x,t)$. 
Fix $(x,t)$ and let $\omega$ be the Riemannian volume vector
at $h(x,t)$ while $\nu$ is the volume vector at $x$. The full
derivative of $h$ is $D [h] : T_x \ensuremath{\mathcal{Q}} \oplus \ensuremath{\mathbb{R}^{}} \to T_{h(x,t)}
\ensuremath{\mathcal{Q}}$. As this is a direct sum of two bundles, we define $D_x [h] :
T_x \ensuremath{\mathcal{Q}} \to T_{h(x,t)}$ by restricting $D[h]$ to the subspace
  $T_x$. We define $\dot{h} = \ensuremath{\frac{\partial {h}}{\partial {t}}}$ in a similar fashion
   except that, by evaluating
at 1, we can say that $\dot{h} \in T_{(h(x,t)}$. Additionally, please note that
$\dot{Q}(h(x,t)) \in T_x\ensuremath{\mathcal{Q}}$.  We shall also take
the liberty of using various equivalences such as $\nabla f \cdot
\ensuremath{\boldsymbol{v}} = D[f] (\ensuremath{\boldsymbol{v}})$ without necessarily mentioning them.


By using the chain rule, we have  
\[
\ensuremath{\frac{\partial {\ensuremath{\rho}}}{\partial {t}}}   = \ensuremath{\lvert {Dh} \rvert} D_h [\ensuremath{\rho}_0] ( \dot{h}) 
+ \ensuremath{\rho}_0(h) \ensuremath{\frac{\partial {[\ensuremath{\lvert {Dh} \rvert}]}}{\partial {t}}}  
\]
and 
\[ \nabla \cdot [ \ensuremath{\rho} \nabla [S] ])   = \ensuremath{\lvert {Dh} \rvert} D_h [\ensuremath{\rho}_0] \circ D_x[h] 
\circ \ensuremath{\boldsymbol{v}}    + \ensuremath{\rho}_0(h) D_x[\ensuremath{\lvert {Dh} \rvert}] \circ \ensuremath{\boldsymbol{v}}   
+ \ensuremath{\rho}  \nabla \cdot \ensuremath{\boldsymbol{v}} . 
\] 


Since $h(Q(x_0,t), t) = x_0$ and $\dot{Q}(x_0,t) = \nabla S (Q)$ 
 by definition , we can eliminate two terms
by noting 

\begin{align*}
\dot{h} (Q, t) + D_Q[h] \circ \dot{Q}  & = 0 \\ 
\dot{h} (Q, t) &= - D_Q[h] \circ \ensuremath{\boldsymbol{v}} (Q)\\
\ensuremath{\lvert {Dh} \rvert} D_h [\ensuremath{\rho}_0] \circ \dot h 
  & = - \ensuremath{\lvert {Dh} \rvert}D_h [\ensuremath{\rho}_0] \circ (  D_x[h] \circ \ensuremath{\boldsymbol{v}} ) \\ 
\end{align*}


Since contractions commute, we have the cancellation of the first terms. 

Replacing $Q$ with $x$ and differentiating, 
we have that 
\[
 D_x[\dot{h}]  
  =  - D_x^2[h] (\ensuremath{\boldsymbol{v}})  - D_x[h] (D_x\ensuremath{\boldsymbol{v}}).
\]
Using the definition of determinant and following the discussion of
the derivative of the determinant, we find that 
\[ \ensuremath{\frac{\partial {[\ensuremath{\lvert {D_x[h]} \rvert}]}}{\partial {t}}} \omega = \sum_i \{D_x[h] \ensuremath{^{\otimes \ensuremath{\mathfrak{n}} \setminus
  i} } \otimes D_x[\dot{h}]\} \nu. 
\]
Combining these two equations, we find, 
\[ \ensuremath{\frac{\partial {[\ensuremath{\lvert {D_x[h]} \rvert}]}}{\partial {t}}} \omega = \sum_i \{D_x[h] \ensuremath{^{\otimes \ensuremath{\mathfrak{n}} \setminus
  i} } \otimes\{ - D_x^2[h] (\ensuremath{\boldsymbol{v}})  - D_x[h] (D_x\ensuremath{\boldsymbol{v}})\}\} \nu. 
\]
Again by a simple computation, we also have, 
\[ D_x[\ensuremath{\lvert { D[h]} \rvert}] (\ensuremath{\boldsymbol{v}}) \omega = \sum_i \{D_x[h] \ensuremath{^{\otimes \ensuremath{\mathfrak{n}}
  \setminus i } } 
\otimes  D^2_x [h] (\ensuremath{\boldsymbol{v}}) \} \nu
\]
which replaces part of the above equation. 

To eliminate the term with the derivative of \ensuremath{\boldsymbol{v}}, one should
recall the definition of the trace of a matrix. Additionally, we shall 
view $D \ensuremath{\boldsymbol{v}}$
 as a matrix as it takes in a tangent vector and
produces a tangent vector. We therefore have  
\[
\begin{split}
 \sum_i \{D_x[h] \ensuremath{^{\otimes \ensuremath{\mathfrak{n}} \setminus
  i} } \otimes\{D_x[h] (D_x\ensuremath{\boldsymbol{v}})\}\} \nu
& =
 D_x [h] \ensuremath{^{\otimes n} } \sum_i  (D\ensuremath{\boldsymbol{v}})^{(i)}(\nu)\\
&  =  D_x[h] \ensuremath{^{\otimes n} } (\ensuremath{\textrm{Tr}} (D\ensuremath{\boldsymbol{v}}) \nu) \\
&= \nabla \cdot \ensuremath{\boldsymbol{v}} \{ D_x[h] \ensuremath{^{\otimes n} } \nu\} \\
&=  \nabla \cdot \ensuremath{\boldsymbol{v}} \{ \ensuremath{\lvert {D [h]} \rvert}\} \omega\}  .
\end{split}
\]


We therefore conclude that 
\[ 
\ensuremath{\frac{\partial {[\ensuremath{\lvert {D_x[h]} \rvert}]}}{\partial {t}}}  + D_x[\ensuremath{\lvert { D[h]} \rvert}](\ensuremath{\boldsymbol{v}}) +  \nabla \cdot \ensuremath{\boldsymbol{v}} \{
 \ensuremath{\lvert {D [h]} \rvert}\} = 0
\]
Thus, all of the terms in the continuity equation do vanish. 


\section{The {H}amilton-{J}acobi equation}

We claim that the Hamilton-Jacobi equation's solution is the action
integral integrated over the trajectories. More precisely, let $V$ be
the potential, $Q$ the set of trajectories, and $S_0$ the initial
action function.  We are assuming that $Q$ is invertible and that
$\dot{Q} (x_0,0) = \nabla S_0 (x_0)$, and $\ddot Q (x_0,t) =
-\nabla V (Q (x_0,t))$. 

Then we claim
that 
\begin{equation} \label{ae:sdef}
S(x,t) = S_0(Q^{-1}(x,t)) + \int_0^t [\frac{1}{2} g(\dot Q, \dot Q) - V
(Q)](Q^{-1}(x, t), s) ds    
\end{equation} 
is the solution to the H-J equation. To establish this, first note
that 
\[
S(Q(x_0,t),t) \ensuremath{=:} \tilde{S} (x_0,t) = S_0(x_0) + \int_0^t
[\frac{1}{2} g(\dot Q, \dot Q) - V 
(Q)](x_0, s) ds    
\] 
satisfies, for fixed $x_0$,    
\begin{equation} \label{ae:lag}
\ensuremath{\frac{d{\tilde{S}}}{d{t}}} (x_0,t) = \ensuremath{\mathcal{L}}(x_0,t) := 
\{\frac{1}{2} g(\dot Q, \dot Q) - V(Q)\}(x_0,t)
\end{equation}
as a trivial application of the fundamental theorem of calculus
establishes.
On the other hand, 
\begin{equation} \label{ae:eng}
\ensuremath{\frac{d{S(Q(x_0,t),t)}}{d{t}}} = \ensuremath{\frac{\partial {S}}{\partial {t}}} (Q, t) + D_Q [S] \circ
\dot{Q}({x_0,t})  .
\end{equation}
The main claim, which requires effort to prove, is that 
\begin{equation} \label{ae:claim} 
\nabla S (Q (x_0,t)) = \dot{Q}(x_0,t)
\end{equation} 
implying 
\begin{equation}\label{ae:qdens}
\nabla S \cdot \nabla S 
= D_Q [S] \circ\dot{Q} d 
= g(\dot{Q}, \dot{Q} ).  l
\end{equation} 
Combining \eqref{ae:qdens}, \eqref{ae:lag}, and \eqref{ae:eng}, we find
\[
\{\ensuremath{\frac{\partial {S}}{\partial {t}}}  + \nabla S \cdot \nabla S 
 - (\frac{1}{2} \nabla S \cdot \nabla S) - V) \}(Q(x_0,t),t)  = 0
\]
which, after replacing $Q$ with $x$, is the Hamilton-Jacobi
equation demonstrating that our claimed solution is 
 a local solution to the
Hamilton-Jacobi equation whenever $Q$ is locally
invertible.\footnote{In Chapter \ref{c:algo}, algorithm A1 requires
  the solution to an equation of the form 
$\ensuremath{\frac{\partial {S}}{\partial {t}}} + \nabla S \cdot \boldsymbol{v} = \frac{1}{2} g(\boldsymbol{v}, \boldsymbol{v})  -
V$. From the above discussion, one can immediately see that
integrating the right hand side over the integral curves of $\boldsymbol{v}$
will yield the solution.}

All that remains to show is that if $S$ is defined by \eqref{ae:sdef}
and that $Q$ is the family of classical trajectories under the
potential $V$ whose initial velocity field is given $\nabla S$, then
\eqref{ae:claim} holds. That is to say, the constraint is
preserved. To prove this, we use the chain rule. 
We start with noting that
\[
D_{x_0} [\tilde{S}] (x_0,t) = D_{x_0} S_0 (x_0) + \int_0^t
\{g( D_{x_0} [\dot Q], \dot Q) - D_Q [V] \circ D_{x_0}[Q]\}(x_0, s) ds     
\] 
and applying the product rule
\[
\ensuremath{\frac{d{g(\dot Q, D_{x_0} [Q] )}}{d{t}}} = g(\ddot{Q}, D_{x_0} [Q])  + g(
\dot Q,  D_{x_0} [\dot Q]) 
\]
 setting us up for an integration by parts leading to 
\[
\begin{split}
D_{x_0} [\tilde{S}] (x_0,t)
=&  D_{x_0} S_0 (x_0) + g(\dot Q, D_{x_0} [Q] )\vert_0^t   
\\&
-\int_0^t 
 {g(\ddot{Q}, D_{x_0} [Q]) -D_Q[V] \circ D_{x_0} [Q]} (x_0, s) ds   
\\ &
=(  D_{x_0} S_0 (x_0) - g(\dot Q (x_0,0) , D_{x_0} [Q] (x_0,0) ) 
+ g(\dot Q, D_{x_0} [Q] ) (x_0,t) 
\\& -  \int_0^t 
 {g(\ddot{Q}, D_{x_0} [Q]) -DV \circ D_{x_0} [Q]} (x_0, s) ds.
\end{split}
\]
As $D_{x_0} [Q] (x_0,0) = \ensuremath{\text{Id}}$, the initial conditions lead to the
cancellation of the first two terms in the above equation. 
The differential of $Q$ can be pulled out to the right in the
integrand leading to
\[
( g(\ddot{Q}, ) -DV )\circ D_{x_0} [Q] (x_0, s) 
= 0 \circ D_{x_0} [Q] (x_0, s)
\]
since $Q_{x_0} (t)$ is a classical trajectory with potential
$V$. The integral's vanishing leaves us with one term.
To finish the claim, note that, evaluating at $(x,t)$,  
\[
D_x [S]  = D_{Q^{-1}} [\tilde S] \circ D_x [Q^{-1}]   
= g(\dot Q, D_{Q^{-1}} [Q] \circ D_x [Q^{-1}]  )  
= g(\dot Q, \ensuremath{\text{Id}}) . 
\]
Taking $x = Q$, this establishes \eqref{ae:claim}. Thus, $S$ is a
solution to the Hamilton-Jacobi equation under the invertibility
assumption of $Q$. 

\section{Semi-classical asymptotics and the Maslov index}

The idea of semi-classical asymptotics in quantum mechanics is quite
well known and we shall be brief. A few references are \cite{Arn89,
  Mas80, dG98} The basic
idea is to take $\hbar \to 0$. Setting it equal to 0 leads to the
classical Hamilton-Jacobi equation. The continuity equation can then
be solved using the classical trajectories. When $\hbar$ is considered
a small
parameter in the problem, this is expected to be a good approximation
to the solution. An immediate problem of this method is that the
classical trajectories are not generically invertible. Thus, the
Hamilton-Jacobi equation does not have a global solution. Undeterred
by this minor inconvenience, one can do some Fourier analysis and,
 by applying the method of stationary phase, we arrive at 
\[
{\ensuremath{\psi}}_n (x, t)  \ensuremath{:=} \sum_{ \{y_j | Q_n(y_j, t) = x \} } \mathcal{R} [Q_n,
 R(\cdot,0)] (y_j,  t)  
e^{i(\mathcal{S}[Q_n, U_{n-1}, S(\cdot, 0)] (y_j, t) -
 \frac{\pi}{2} \mu_j)}.
\]
In words, we are assuming that in a neighborhood of each point, the
function $Q$ is locally invertible. We look at the inverse image and
generate the $S$ along the trajectories using the action formula. We
append the Jacobian correction factor and do a sum over the terms. It
is important that the terms have the right phase which is accounted
for by
the Maslov index $\mu$ over the trajectory. Although it is a very
important symplectic geometric object associated with Lagrangian
submanifolds,  we shall view it in a very
simple-minded way. The factor is just the number of zeroes of the
determinant of $Q$ that the trajectory crosses. The zeroes correspond
to caustics and these are the places where invertibility of $Q$ begins
to fail. The reason this factor arises from our perspective is
clear. Remove the absolute value signs around the determinant and the
square root will then pick up factors of $i$ (or $-i$)  as a
trajectory goes through a zero. Noticing the multiplicity arises by
considering a slight perturbation of the surrounding flow and showing
that the overall phase factor is unaffected by the perturbation.

\clearpage
\bibliographystyle{abbrv}
\bibliography{references}
\end{document}
